\documentclass[11pt,a4paper]{article}

\usepackage[margin=24mm]{geometry}
\usepackage[T1]{fontenc}
\usepackage[utf8]{inputenc}
\usepackage{lmodern}
\usepackage{microtype}
\usepackage{amsmath}
\usepackage{xcolor}
\usepackage{hyperref}
\usepackage{listings}

\hypersetup{
  colorlinks=true,
  linkcolor=blue!55!black,
  urlcolor=blue!55!black,
  pdftitle={Manual for compare_urdf_dh_frames.m}
}

\definecolor{codebg}{RGB}{248,248,248}
\definecolor{codeframe}{RGB}{210,210,210}
\lstset{
  basicstyle=\ttfamily\small,
  backgroundcolor=\color{codebg},
  frame=single,
  rulecolor=\color{codeframe},
  breaklines=true,
  columns=fullflexible,
  keepspaces=true,
  showstringspaces=false
}

\title{Concise Manual for \path{compare_urdf_dh_frames.m}}
\author{Robot System Identification utilities}
\date{\today}

\begin{document}
\maketitle

\section{Purpose}

\path{compare_urdf_dh_frames.m} checks a standard Denavit--Hartenberg (DH)
model generated from a URDF serial chain. At a user-supplied joint
configuration, it reconstructs the DH frames, evaluates the same configuration
with the URDF model, and compares:
\begin{itemize}
  \item every URDF joint axis with the corresponding DH axis $z_{i-1}$; and
  \item the URDF and DH reconstructions of the selected tool pose.
\end{itemize}
It also creates a figure containing the URDF model, the DH chain, and an
aligned overlay. The comparison is diagnostic and does not modify its inputs.

\section{Requirements and inputs}

The function requires MATLAB with Robotics System Toolbox. Add the utility
folder to the MATLAB path before calling it. Its three file inputs are:
\begin{description}
  \item[URDF file] The source robot description.
  \item[DH CSV] The generated \path{dh_table.csv}, including the columns
    \path{joint}, \path{a_m}, \path{alpha_rad}, \path{d_m},
    \path{theta_rad}, \path{t}, and \path{s}.
  \item[Kinematics JSON] The generated \path{urdf_kinematics.json}, containing
    the base-to-first-DH-frame and last-DH-frame-to-tool transforms.
\end{description}
These generated files are documented in
\href{run:urdf_to_standard_dh_manual.pdf}{\path{urdf_to_standard_dh_manual.pdf}}.

The vector \path{q} must contain one finite joint coordinate per CSV row, in
the same order as the CSV. Angles are expressed in radians and prismatic
coordinates in metres.

\section{Function interface}

\begin{lstlisting}[language=Matlab]
result = compare_urdf_dh_frames( ...
    urdfPath, dhCsvPath, kinematicsJsonPath, q, ...
    'BaseLink', baseLink, ...
    'ToolLink', toolLink, ...
    'FrameScale', 0.08, ...
    'ShowVisuals', true);
\end{lstlisting}

\path{BaseLink} and \path{ToolLink} are required. They select the URDF chain
to compare and prevent unrelated branches or fixed frames from entering the
test. \path{FrameScale} is the plotted coordinate-axis length in metres;
its default is $0.08$. \path{ShowVisuals} enables URDF visual meshes and is
\path{true} by default. Setting it to \path{false} can make complex models
faster and clearer to inspect.

\section{Example}

\begin{lstlisting}[language=Matlab]
addpath(fullfile('src', 'utils', 'UrdfToDH'));

robotRoot = fullfile('src', 'Identification_Luna_DB3.5', ...
    'RobotDescription');
q = zeros(1, 7);

result = compare_urdf_dh_frames( ...
    fullfile(robotRoot, 'urdf', 'arm_with_None.urdf'), ...
    fullfile(robotRoot, 'DH', 'dh_table.csv'), ...
    fullfile(robotRoot, 'DH', 'urdf_kinematics.json'), ...
    q, ...
    'BaseLink', 'starman_arm_link_cart', ...
    'ToolLink', 'starman_arm_ids');
\end{lstlisting}

Run the comparison at several representative configurations, not only at the
zero pose. Configurations must respect the URDF joint-coordinate convention
used to generate the DH model.

\section{Computation and displayed result}

For row $i$, the function evaluates the standard-DH transform
\begin{equation}
 A_i = \operatorname{RotZ}(\theta_i^*)
       \operatorname{TransZ}(d_i^*)
       \operatorname{TransX}(a_i)
       \operatorname{RotX}(\alpha_i),
\end{equation}
where
\begin{align}
 \theta_i^* &= \theta_i + (1-t_i)s_iq_i, &
 d_i^* &= d_i + t_is_iq_i.
\end{align}
The residual transforms from the JSON file align the reconstructed DH chain
with the selected URDF base and tool.

The figure has three panels: URDF model and body frames, generated DH frames,
and their overlay. DH frames are labelled \path{D0}, \path{D1}, and so on;
URDF joint axes are labelled \path{UJ1}, \path{UJ2}, and so on. The command
window reports the tool position and orientation errors and, for each joint,
the axis-angle error and shortest distance between the URDF and DH axis lines.

\section{Returned structure and interpretation}

The returned \path{result} structure contains joint names and configuration,
the URDF and DH tool transforms, every reconstructed DH frame, joint-axis
angle and line-distance errors, tool position and orientation errors, and the
figure handle.

For a consistent conversion, all reported errors should be close to numerical
zero, subject to the precision of the URDF and generated files. A large tool
error indicates a kinematic mismatch. A large per-joint angle error indicates
an axis-direction mismatch, while a large line distance indicates that the
two compared axes are displaced.

Under the standard-DH convention, URDF joint $i$ is associated with
$z_{i-1}$, not $z_i$. URDF body frames and DH frames are defined for different
purposes and therefore do not need to coincide visually; joint-axis and final
tool-pose agreement are the relevant checks.

\end{document}
